\section{Fundamentos de la muestreo}

Antes de discutir cómo aprender distribuciones de datos complejas, el instructor repasa una pregunta fundamental: \textbf{si ya conocemos una distribución de probabilidad, ¿cómo se extraen muestras de ella?}

\subsection{Muestreo por transformación inversa (Inverse Transform Sampling)}

\subsubsection{Caso discreto: Muestreo de una distribución Categorical}

Consideremos una distribución categorical con $K$ estados, donde cada estado $k$ tiene probabilidad $p_k$, y $\sum_{k=1}^K p_k = 1$. El método clásico para extraer muestras de esta distribución es el \textbf{muestreo por transformación inversa}:

\begin{enumerate}
    \item Calcular la función de distribución acumulada (CDF): $F(k) = \sum_{i=1}^{k} p_i$
    \item Extraer un valor $u$ de una distribución uniforme $U \sim \text{Uniform}(0, 1)$
    \item Encontrar el estado $k$ que satisface $F(k-1) < u \leq F(k)$; dicho estado es el resultado del muestreo
\end{enumerate}

Intuitivamente, los estados con mayor probabilidad ocupan un ``intervalo'' más ancho en el eje $y$ de la CDF, por lo que la probabilidad de que una muestra uniforme caiga dentro de ese intervalo también es mayor.

\subsubsection{Caso continuo}

Para una variable aleatoria continua $X$, cuya función de densidad de probabilidad (PDF) es $f(x)$ y cuya CDF es:

$$F(x) = \int_{-\infty}^{x} f(t) \, dt$$

Los pasos del muestreo por transformación inversa son:

\begin{enumerate}
    \item Extraer $u \sim \text{Uniform}(0, 1)$
    \item Calcular $x = F^{-1}(u)$
\end{enumerate}

\begin{importantbox}{Fórmula clave del muestreo por transformación inversa}
Sea $F$ la CDF de una distribución continua y $F^{-1}$ su función inversa, entonces:
$$X = F^{-1}(U), \quad U \sim \text{Uniform}(0,1) \implies X \sim F$$
Suposición clave: \textbf{es posible calcular la función inversa de la CDF, $F^{-1}$}.
\end{importantbox}

\subsubsection{Ejercicio de clase: Muestreo de $f(x) = \frac{3}{8}x^2$}

El instructor propone un pequeño ejercicio: dado el PDF $f(x) = \frac{3}{8}x^2$, con $x \in {[}-2, 2{]}$, derive la fórmula de muestreo por transformación inversa.

Solución:
\begin{enumerate}
    \item Calcular la CDF: $F(x) = \int_{-2}^{x} \frac{3}{8}t^2 \, dt = \frac{1}{8}(x^3 + 8)$
    \item Plantear $u = F(x)$ y despejar $x$: $u = \frac{1}{8}(x^3 + 8) \implies x = (8u - 8)^{1/3} = 2(u-1)^{1/3}$
    \item Corrección, derivación más cuidadosa: $x^3 = 8u - 8 \implies x = \sqrt[3]{8u - 8}$
\end{enumerate}

Por lo tanto, el proceso de muestreo consiste en: primero extraer $u \sim \text{Uniform}(0,1)$ y luego calcular $x = (8u - 8)^{1/3}$.

\subsection{Muestreo de la distribución gaussiana: Transformación Box-Muller}

\subsubsection{Planteamiento del problema}

El PDF de la distribución normal estándar $\mathcal{N}(0, 1)$ es:

$$f(x) = \frac{1}{\sqrt{2\pi}} e^{-x^2/2}$$

Su CDF $\Phi(x) = \int_{-\infty}^{x} f(t) \, dt$ \textbf{no tiene una función inversa en forma analítica}, por lo que no se puede utilizar directamente el muestreo por transformación inversa.

\begin{warningbox}{Trampa de la CDF gaussiana irreversible}
Aunque el PDF de la distribución gaussiana tiene una forma analítica elegante, su función inversa (función probit) no existe en forma cerrada. Esto significa que no se puede aplicar directamente el método más simple de muestreo por transformación inversa para distribuciones gaussianas. En la práctica es necesario recurrir a trucos como sustitución de variables.
\end{warningbox}

\subsubsection{Derivación de la Transformación Box-Muller}

La idea central de la Transformación Box-Muller es: \textbf{considerar simultáneamente dos muestras independientes de distribución normal estándar, aplicando una transformación a coordenadas polares para que cada componente pueda muestrearse mediante transformación inversa}.

Sea $X, Y \overset{\text{i.i.d.}}{\sim} \mathcal{N}(0, 1)$, cuya densidad conjunta es:

$$f_{X,Y}(x, y) = \frac{1}{2\pi} e^{-(x^2 + y^2)/2}$$

Introduciendo variables en coordenadas polares $D = X^2 + Y^2$ (radio al cuadrado) y $\Theta$ (ángulo), se puede demostrar:

\begin{itemize}
    \item $\Theta \sim \text{Uniform}(0, 2\pi)$ (distribución uniforme del ángulo, que se extrae directamente)
    \item $D \sim \text{Exponential}(1/2)$, es decir, $f_D(d) = \frac{1}{2}e^{-d/2}$ (distribución $\chi^2$ con 2 grados de libertad)
\end{itemize}

Para la distribución de $D$, la CDF es $F_D(d) = 1 - e^{-d/2}$, cuya función inversa es:

$$d = F_D^{-1}(u) = -2\ln(1 - u)$$

Dado que $1-u$ y $u$ son ambos $\text{Uniform}(0,1)$, esto se simplifica a $d = -2\ln(u)$.

\begin{importantbox}{Fórmula completa de la Transformación Box-Muller}
Dadas dos muestras independientes de distribución uniforme $U_1, U_2 \sim \text{Uniform}(0,1)$:
$$X = \sqrt{-2\ln U_1} \cdot \cos(2\pi U_2)$$
$$Y = \sqrt{-2\ln U_1} \cdot \sin(2\pi U_2)$$
Entonces $X$ y $Y$ son dos muestras independientes de distribución normal estándar. Una sola transformación genera simultáneamente dos muestras gaussianas.
\end{importantbox}

\subsection{Muestreo de distribuciones gaussianas multivariantes y Reparameterization}

\subsubsection{Distribución gaussiana isótropa}

Para una distribución normal multivariante estándar $\mathcal{N}(\mathbf{0}, \mathbf{I})$, dado que las dimensiones son independientes, se puede aplicar la Transformación Box-Muller a cada dimensión por separado.

\subsubsection{Distribución gaussiana anisótropa}

Para una distribución gaussiana multivariante general $\mathcal{N}(\boldsymbol{\mu}, \boldsymbol{\Sigma})$, el método de muestreo es la \textbf{reparametrización}:

$$\mathbf{x} = \boldsymbol{\mu} + \mathbf{L} \boldsymbol{\epsilon}, \quad \boldsymbol{\epsilon} \sim \mathcal{N}(\mathbf{0}, \mathbf{I})$$

Donde $\mathbf{L}$ satisface $\boldsymbol{\Sigma} = \mathbf{L}\mathbf{L}^T$ (por ejemplo, descomposición de Cholesky).

\begin{importantbox}{Truco de reparametrización}
Muestrear de $\mathcal{N}(\boldsymbol{\mu}, \boldsymbol{\Sigma})$ es equivalente a:
$$\mathbf{x} = \boldsymbol{\mu} + \boldsymbol{\Sigma}^{1/2} \boldsymbol{\epsilon}, \quad \boldsymbol{\epsilon} \sim \mathcal{N}(\mathbf{0}, \mathbf{I})$$
Este truco no solo es fundamental para el muestreo, sino clave para implementar la propagación del gradiente en VAEs: aísla la aleatoriedad en $\boldsymbol{\epsilon}$, haciendo que el proceso de muestreo sea diferenciable respecto a $\boldsymbol{\mu}$ y $\boldsymbol{\Sigma}$.
\end{importantbox}

\begin{knowledgebox}{¿Por qué los modelos generadores prefieren las distribuciones gaussianas?}
Las distribuciones gaussianas son omnipresentes en los modelos generativos debido a: (1) conocemos cómo muestrearlas de manera eficiente; (2) poseen excelentes propiedades matemáticas (como que la divergencia KL tiene forma cerrada); (3) el teorema del límite central garantiza que muchas distribuciones tienden asintóticamente a la gaussiana; y (4) bajo restricciones de media y varianza, la distribución gaussiana maximiza la entropía.
\end{knowledgebox}

\subsection{Resumen de la sección}

Cuando la forma de una distribución de probabilidad es conocida, podemos realizar muestreo mediante transformación inversa, sustitución de variables (como la Transformación Box-Muller) y reparametrización. Sin embargo, las distribuciones de datos reales (como la de imágenes naturales) tienen formas desconocidas y extremadamente complejas; esto es precisamente lo que motiva el uso de modelos generativos profundos.

\newpage
